\chapter{Methods} % (fold)
\label{chap:methods}

Following our discussion on the nature of chirps and their detection methods,
this chapter outlines the pipeline developed for chirp detection. This pipeline
comprises four key components: 

\begin{enumerate}
  \item \textbf{Feature extraction from grid recordings:} This 
        initial step involves generating spectrograms to identify chirps.
  \item \textbf{Chirp detection:} In this phase, chirps are detected on the
        spectrograms, undergoing various transformations and filters.
  \item \textbf{Feature extraction for chirp assignment}: This involves
        analyzing individual chirps on the spectrograms to extract features
        necessary for assigning them to specific fish.
  \item \textbf{Chirp assignment}: The final step reduces the assignment
        problem into a binary classification task, addressed with another
        neural network.
\end{enumerate}

The pipeline utilizes two separate feature extraction processes and two neural
networks to detect and assign chirps. Training and evaluating these networks
are important steps in the process. Therefore, before jumping into the
pipeline's specifics, this chapter will first detail the training process for
each neural network, including the data used for training.

\section{Chirp detection}
\label{sec:cpd}

The issue of chirp detection is where most of the complexity of the pipeline
lies and most of the training data had to be generated. The following sections
will explore the different strategies I used to identify chirps in a set of
simple recordings and how I used them to build a training dataset for complex
grid datasets.

\subsection{Datasets}
\label{sec:rawdata}

I used multiple datasets to extract training data for the object detection
model. The first dataset was recorded during competition experiments by
\textcite{raabElectrocommunicationSignalsIndicate2021} in \textit{A.
leptorhynchus}. They staged competitions of two fish over access to a shelter
in experimental tanks and recorded electric signals on an electrode grid as
well as spatio-temporal behaviors using an near-infrared camera. The recordings
were made with a grid of electrodes and fish were able to move freely in the
tank. This provided to be vital for generating a diverse training dataset:
While fish were competing, they produced a large number of chirps with a much
higher diversity in their characteristics as compared to stationary
chirp-chamber recordings. Additionally, the movement of the fish resulted in
much more challenging spectrograms. This illustrated early on how diverse the
chirps of \textit{A. leptorhynchus} can be and how movement increases the
complexity of chirp detection. The disadvantage of this setup is that if I
wanted to use this data to generate training data, computed spectrograms would
have needed to be annotated by hand to extract the chirps. And in an ideal
world, a dataset used to train an object detection model would contain multiple
thousands of images, each with multiple instances of the object of interest.
This is why I decided to add a second available dataset of chirp-chamber
recordings.\label{sec:ccdata}

The chirp-chamber recordings were produced from single individuals of
\textit{A. leptorhynchus} that were restrained to a PVC tube using mesh on both
ends, as described in \textcite{stocklEncodingSocialSignals2014}. The fish
were stimulated by a pair of electrodes that were placed on the left and right
side of the chamber. The two recording electrodes where positioned at the head
and tail of the fish. Because the stimulation electrodes are isopotential to
the field of the fish, the stimulus did not interfere with the EOD and hence
was not recorded on the recording electrodes. The fish were then stimulated by
a sine-wave mimicking the presence of a conspecific. The frequency of the sine-wave was varied between $\pm$ 30 \si{\hertz} around the EOD$f$ of the fish.
While this dataset included a large number of chirps, the diversity of chirps
was limited to mostly very short type 1 and type 2 chirps
\parencite{englerSpontaneousModulationsElectric2000a}.

Because of the lack of diversity in this dataset, I also included a third
dataset of various experiments that mostly included two fish interacting. The
dataset was recorded by \textcite{obotiWhyBrownGhost2023} and most fish were
spatially restricted to a certain degree, meaning that the strongest signal on
one of two recording electrodes was always the EOD of the fish of interest.
This dataset includes a large variety of different experiments, all of which
are documented in the original publication.

\subsection{Training data acquisition}
\label{sec:dettraindata}

The first step in training the object detection model was to generate a large
number of training images, i.e., spectrograms with chirps. As the datasets that
I had access to at the beginning of the project were very large, converting all
the data to spectrograms and annotating the spectrograms by hand was not
feasible. This is why I decided to start by modeling chirps in simulated grid
recordings. At least in theory, this provides the ability to generate an
infinite amount of training data without the need of manual annotation.

\subsubsection{Initial training on modeled chirps}

To train an object detection model on spectrograms, merely modeling chirps is
not enough. To get a realistic spectrogram, a simulated recording of multiple,
moving and communicating fish is necessary. In this section, I will describe
the process of generating such a simulated recording and how to extract
training data from it. The first step was to model the frequency excursions
that happen during chirps.

\emph{Simulating a chirp.} To model chirps, I first needed to understand and
parameterize the chirps that are present in the recordings. As the only dataset
that I had access to at the time was the competition dataset of
\textcite{raabElectrocommunicationSignalsIndicate2021}, I decided to use the
chirps that were present in this dataset to model chirps. As a first step, I
had to extract instantaneous frequency estimates of the EOD of a chirping fish.
To do this, I used the chirps that were already sparsely detected by the
previous chirp detection approaches as first estimates on where to look for
more. I only considered the chirps of fish with the highest baseline EOD$f$.
The chirps of the fish with the lower baseline EOD$f$ would cross the frequency
band of the upper fish, rendering them unusable. I used a set of heuristics to
extract the instantaneous frequency for each chirp, which are described in
detail in the \hyperref[appendix]{appendix}.

To model chirps, I then fitted a Gaussian frequency modulation to the centered
instantaneous frequency of each chirp. The Gaussian modulation had a height
$h$, a width $w$ and a kurtosis $k$. The standard deviation was then computed
from the width and kurtosis:

\begin{equation}
  g(t) = h \exp\left(-\frac{1}{2}\left(\frac{t - \mu}{\sigma}\right)^{2}\right)^{k} \quad \text{with} \quad
  \sigma = \frac{w / 2}{2 \log(10)^{k/2}}
\end{equation}

Where $t$ is the time, $\mu$ is the center of the chirp and $k$ is the
kurtosis. A kurtosis of 1 corresponds to a perfect Gaussian. In addition to
this rather simple model, I also experimented with mixtures of up to three
Gaussian modulations. To achieve realistic type 1 and type 2 chirps, at least
two Gaussian curves should be needed as these chirps often contain an
undershoot or down modulation of the frequency at the end. The third Gaussian
was introduced to allow facilitate degree of freedom. But as model complexity
increases, so do the fitting errors, which is why I decided to stick with the
simple Gaussian model. After fitting I then linearly interpolated the three
parameters of the monophasic model. I sorted all parameter groups by the width
of the Gaussian and then interpolated all parameters to fill the gaps in the
parameter space. This allowed modeling chirps of any duration, height and
kurtosis, limited by the minima and maxima of the real chirps in the dataset.

\emph{Simulating a chirping fish.}\label{sec:simulatingachirpingfish} Moving
from frequency space into the time domain, the next step involves simulating
the EOD of \textit{A. leptorhynchus}, containing the simulated chirps. The EOD,
a periodic signal with a fundamental frequency and harmonics, can be modeled as
a sum of sinusoids, if its Fourier spectrum is known. Utilizing the
\href{https://github.com/janscience/thunderfish/}{\texttt{thunderfish}}
package, we can easily model an EOD, including specific amplitudes and phases
for different taxa.

To simulate an EOD's evolution over time, I start by selecting a baseline
frequency between 600 and \SI{1100}{Hz} from a uniform distribution, reflecting
the typical range for \textit{A. leptorhynchus}. To introduce chirps,
parameters are randomly drawn from the previously described interpolated
Gaussian model and  the resulting frequency curve is added to the EOD$f$
vector. The simulation randomly determines the number of chirps, ensuring no
more than one chirp per second and a minimum interval of
\SI{200}{\milli\second} between chirps. To achieve this, a random number of
chirps (constrained by the maximum rate) is randomly positioned onto the full
frequency trace and chirps that are too close are discarded.  Notably, chirps
also involve a change in EOD amplitude, typically a slight reduction. This
effect is simulated by inverting the Gaussian used for frequency modulation,
adjusting its scale, and multiplying it with the EOD. The resulting amplitude
is randomly set between 100\% and 20\% of the EOD's maximum amplitude.
Reconstructing the EOD as a sum of sine waves using the EOD$f$ vector and
additional parameters for phases and amplitudes results in a semi-realistic
simulation of a chirping fish that oscillates between negative one and one.

At this point, the simulation resembles a recording one would get from a
chirp-chamber setup: A single fish that is not moving and producing chirps,
recorded on two electrodes positioned at its head and tail. This is not enough
to emulate the complexity of a grid recording of multiple, moving fish. So as a
next step, let's add movement to the artificial fish.

\emph{Simulating a moving fish.} An electrode grid reduces the spatial
information to a two-dimensional plane. Because the power of the EOD decays
with distance, it is informative about the distance of the fish to the
recording electrode. In such electrode grid recordings, the positions of fish
in two dimensional space are reflected by the distribution of the EOD
amplitudes on the electrodes. And when a fish moves, the amplitude of the EOD
on each electrode changes. To simulate this, I first simulated the
\enquote{true} movement of a fish over time, from which I then calculated the
distances to all electrodes and the corresponding \acrfull{AM}. As the exact
movement patterns do not matter for this simulation, I used a simple random
walk. The algorithm gets a starting position, a target duration and two
\acrfullpl{PDF}, one for the direction and one for the distance of a single
discrete step. Manipulating the probability densities as well as the number of
steps per time strongly influences how the simulated fish moves. The PDF of the
step distance was a Gamma distribution normalized to the number of steps per
second with a peak at \SI{0.2}{\meter\per\second}, which is equivalent to
velocities observed in the field. The direction PDF was a sum of two Gaussian
distributions with a peak at \SI{0}{\degree} and \SI{180}{\degree}. Increasing
the width of the Gaussians for the direction PDF would make the movement
trajectories much less smooth, as the fish would change direction more
frequently. And \emph{vice versa}, narrowing both Gaussians would make them
smoother. Increasing the weighting of the forward Gaussian, i.e., the one at
\SI{0}{\degree} relative to the backward Gaussian, would decrease the number of
times the fish moves backwards. The forward Gaussian was set to a standard
deviation of 0.2 and 0.1 for the backward Gaussian on an angular unit circle.
The PDF of the backward Gaussian was weighted as 20\% of that of the forward
Gaussian. The chosen parameters ought to simulate the characteristic
forward-backward movement of knife fish like \textit{A. leptorhynchus}. As a
result, we get a vector that describes the position of a fish over time in a
two-dimensional plane. Since the random walk often breached the boundaries of
simulated space, I included boundaries that the fish could not cross. The
resulting positions were \enquote{folded} back into the simulated space (figure
\ref{fig:space_folding}) after generating the random walk.

\begin{figure}
  \centering
  \includegraphics[width=\textwidth]{../figures/space_folding.pdf}
  \caption{Simulated fish movements for three fish. The circles indicate
  electrodes and the black border is the boundary of the virtual space. The
  positions are simulated by a random walker that is capable of stepping into
  any direction of a circle around it. The trajectory of each step underlies a
  bimodal PDF as real fish mostly move forwards or backwards. The
  distance of each step is drawn from a Gamma distribution limited by
  velocities estimated from fish in their natural habitat with a peak at
  \SI{0.2}{\meter/\second}. To maximize computational efficiency, the positions
  are first simulated and then folded back across the boundaries of the virtual
  space until they are contained in the boundaries.}
  \label{fig:space_folding}
  \vspace{0.5cm}
\end{figure}

From these $x$ and $y$ positions, the next step was to simulate the power of
the EOD on each electrode depending on the position of the fish. For this, I
defined a grid of 16 virtual electrodes spaced \SI{50}{\centi\meter} apart in
the simulated space (figure \ref{fig:space_folding}). In the previous steps, we
covered everything needed to simulate a fish that is recorded on a single
electrode. To simulate how this fish would be recorded on the 16 virtual
electrodes, we can simply duplicate the signal 16 times. We then
retrospectively adjust the AM so that it is proportional to the distance between the fish and each of the 16 electrodes. By using Pythagoras' theorem
and the known electrode positions as well as the known fish positions, we can
calculate, for each point in time, the distance between the virtual fish and
each virtual electrode. As the amplitude of the EOD decays with the squared distance, inverting, squaring and normalizing the distance gives us the
amplitude of the EOD$f$ on each electrode. Normalization is necessary, as the
simulated EOD is restricted between 1 and $-1$, while the distance
can become much greater. To obtain an amplitude modulation that we can add to
the EOD, this AM computed from the distances also has to be restricted to
between 1 and $-1$. The main assumption here is that the power of
the electric field decays equally in all directions around the fish. In
reality, the field of the fish is an oscillating dipole with two lobes. But
because I assumed, that this property is not reflected in the appearance of the
sum spectrogram across all recording channels, I ignored it for now. Adding
this distance transformed into an AM to every single EOD gives us a matrix that
contains the EOD including position-dependent AMs. The result is a simulated
electrode grid recording of a single, moving fish that produces chirps.

\emph{Simulating multiple fish.} A single fish is not enough to emulate the
complexity of a grid recording in the field. In reality, the frequencies of
fish often overlap or cross each other, not to speak of the chirps that are
produced. To simulate this, I simply repeated the process of simulating a
single fish a random number of times (between one and five fish
with a minimum $\Delta\text{EOD}f$ between two fish of \SI{20}{\hertz}).
Combining the separate fish is then simply the sum of the matrices that contain
the EOD over time on all electrodes. The result is a simulated grid recording
of multiple, moving fish that produce chirps.

\emph{Simulating noise.} The simulated recordings are very clean and do not
resemble the complexity of a real recording. I added noise at different stages
of the simulation pipeline. Before doing so, we first have to consider how this
translates to spectrograms as this is the only feature that the model is able
to consume. Adding noise to the EOD$f$ would make the frequency band of a
single fish on a spectrogram broader. Adding noise to the EOD however, would
change the \acrfull{SNR}, i.e., the contrast between the \enquote{background}
and the frequency band of a single fish on a spectrogram. These sources of
noise can be easily implemented, given that they are static over time. So as a
first step, I added band-limited Gaussian noise to the baseline EOD$f$ of each
fish (cutoff at 0.0 \si{\hertz} and 0.01 \si{\hertz}, $\sigma = 5
\si{\hertz}$). With the frequency trace of each chirp, I multiplied
band-limited Gaussian noise (cutoff at 600 \si{\hertz}, $\sigma = 1
\si{\hertz}$) as the noisiness of the instantaneous frequency appeared to
increase proportionally with the increase in frequency relative to the baseline
EOD$f$. To the waveform of an EOD I additionally added Gaussian noise ($\sigma
= 0.01$) to simulate its presence in the recordings of single fish. The part
where it gets complicated is the noise that is non-stationary, meaning that it
substantially changes over time. This kind of noise is particularly strong in
laboratory recordings as many appliances, like lights, cables, and other
equipment, introduce electromagnetic noise. Due to its heterogeneity, it is
impossible to simulate \textit{ad-hoc}. It depends on the building, the room,
the time of day and many other factors. So instead of simulating it, I used the
noise floor from real recordings and added it to the simulated recordings. To
achieve this, I selected areas from the Raab dataset which were known for
either contain little or no chirps. I scaled the simulated recording so that it
fits the units of the recorded snippet and added it to my simulated ensemble of
fish to get a realistically noisy backdrop. This still contains the EODs of the
real fish from the Raab dataset, but as they do not produce any chirps, they
can be considered as additional noise.

The result is an artificial grid recording that contains multiple, moving fish
producing chirps and that has a realistic noise floor. As I have simulated
them, I also know when which fish produced a chirp and what the chirp looked
like. But this is not yet a dataset that can be used to train an object
detection model. To train a detection model, we need spectrograms that contain
the frequency bands of chirping fish and labels that contain the coordinates of
bounding boxes around the chirps. So now is a good time to introduce the format
of the training dataset and how spectrograms are computed.

\emph{Spectrogram conversion.}\label{sec:spectrogramconversion} Computing
spectrograms is a computationally expensive task, especially with weeks of
recordings. The two factors that limit how much data of one recording should be
converted to spectrograms at once are the available memory and the time it
takes to compute. Using modern graphics cards can decrease computational time
but limits the amount of data in terms of memory, as \acrshort{GPU}s usually
have less memory than \acrshort{CPU}s. Using an NVIDIA GeForce RTX 4080 with 16
GB of memory, I was able to convert three minutes of recordings from 16
channels sampled at 20 \si{\kilo\hertz} at once (nfft=4096, hop length=410). Converting three minutes of recording takes
approximately 300 \si{\milli\second}. All channels in a grid recording are
converted to a sum spectrogram. Additionally, the spectrogram is transformed
from the power scale to the decibel scale. The windows that I extracted from
the recordings always had an overlap of one second. Another issue is that the
resulting spectrograms are quite large: Finding an event of a few tens of
milliseconds in three minutes of data is like finding a needle in a haystack.
Additionally, the relevant part of the frequency spectrum only reaches up to 2
\si{\kilo\hertz} of the available 10 \si{\kilo\hertz}. To make this data easier
to process for the neural network, I split the spectrograms into 15
\si{\second} chunks with 1 \si{\second} overlap. I also split them along the
frequency dimension into the ranges 0-1000 \si{\hertz}, 500-1500 \si{\hertz},
1000-2000 \si{\hertz}. To generate training data, each of these spectrogram
chunks is saved to disk as a greyscale PNG image.

\emph{Generating labels.} To train a model for predicting the position of
chirps on spectrograms, we need labels that allow the model to learn. A common
way to label the position of objects in images, or in this case, chirps on
spectrograms, is, to use bounding boxes. A bounding box is a rectangular box
around an object of interest that is defined by the coordinates of the four
corners. In the case of simulated data, these labels can be generated from the
known chirp times, widths and heights. To do this, I had to find a
transformation that approximates the width and height of a chirp on a
spectrogram given the same on the instantaneous frequency. As we have seen in
figure \ref{fig:multifish_chirps}, chirps look different on spectrograms
compared to their instantaneous frequency trace. I approximated the width of
the chirp on the spectrogram by adding a term to the original width that is
dependent on the sample rate $f_s$ and spectrogram: $\text{nfft} / f_{s} \times
0.8$. This is necessary because the minimum possible width of a chirp is the
width of the window that is used to compute the spectrogram. The upper boundary
of the bounding box was then set to the maximum frequency of the chirp and a
certain frequency padding given by $f_{s}/\text{nfft}\times5$. The lower
boundary was set to the minimum frequency of the chirp minus the same frequency
padding.

\emph{Generating the dataset.} The dataset is then generated by converting the
simulated grid recordings to spectrograms and saving them to disk as PNG images
and generating labels from the known chirp times, widths and heights. The
labels for each spectrogram were saved as \texttt{txt} files with the same
file name as the image file and coordinates transformed relative to the image
dimensions. This is the common dataset format for YOLO-style object detectors.
The result was the first dataset that I used to train the object detection
model. From now on, I will refer to this dataset as the \emph{simulated
training dataset}.

\subsubsection{Manual annotation with a human in the loop}

In addition to simulated data, real data was also used to train the object
detection model. To efficiently manage the annotation workload, subsets known
to contain chirps were extracted from the Raab datasets. These subsets were
pre-annotated using a faster R-CNN model, an established CNN for object
detection that was fine-tuned on simulated data
\parencite{renFasterRCNNRealTime2016}. The pre-annotations were imported into
\href{https://labelstud.io/}{\texttt{label-studio}}, a web-based, open-source
annotation tool designed for human-in-the-loop annotation processes
\parencite{tkachenkoLabelStudioData2022}. This tool facilitated the refinement
of pre-annotations and the addition of new ones by human annotators. The
refined annotations, from now on termed the \emph{recorded training dataset},
were exported and progressively expanded through cycles of model training,
pre-annotation, human correction, and export. The final dataset comprised 537
images of spectrograms with 14,832 annotations, averaging 27.6 chirps per
image.

\subsubsection{Hybrid simulations}

The last and probably most important dataset that I used to train the object
detection model is what I will refer to as the \emph{hybrid training dataset}.
This is a combination of a fully simulated grid recording with the difference
that the instantaneous frequency evolution of each chirp is not generated by a
simple Gaussian model as it was previously, but extracted from real recordings.
Previously, I introduced the \hyperref[sec:ccdata]{chirp-chamber dataset} which
has the advantage that the fish are solitary and stationary, facilitating the
extraction of the instantaneous frequency. Additionally, these recordings are
differential, meaning that they contain just a single channel signal. Chirps
were simply detected as peaks on the instantaneous frequency. I used a set of
hand-tuned heuristics to extract chirps from these chirp-chamber recordings
that are described in detail in the \hyperref[appendix]{appendix}. The number
of peaks discarded in each filtering step are visualized in figure
\ref{fig:chirp_baseline_smoothing} A.

This resulted in 9,082 extracted chirps. The resulting chirps still had an
often jittery baseline frequency. If I would put these chirps directly into the
simulation pipeline, the area of the spectrogram around the chirp would be
different from the rest of the baseline EOD$f$ band. This is a feature a neural
network could learn but that is not present in real data. To prevent this, I
had to find a way to make the baseline frequency, i.e., everything left and
right of the chirp, as smooth as possible while keeping the natural variability
of the EOD$f$ during a chirp.

\begin{figure}
  \centering
  \includegraphics[width=\textwidth]{../figures/plot_chirp_baseline_crossfade_smoothing.pdf}
  \caption{
    Visualization of the filter contributions and frequency post-processing
    steps. (A) Shows how many of the detected peaks in the EOD$f$ were
    discarded by each filter threshold. The thresholds were chosen to remove
    false positives and to keep the chirps that are most likely to be real. For
    details see \hyperref[appendix]{appendix}. (B) Illustrates how the baseline
    instantaneous frequency around a chirp was smoothed. I used the smoothed
    baseline to find the beginning and end of the chirp and applied a
    cross-fade to the instantaneous frequency to make the baseline around the
    chirp as smooth as possible while retaining the natural variability of the
    instantaneous frequency during a chirp. The result is a centered
    instantaneous frequency that can be used to simulate a chirping fish.
  }
  \label{fig:chirp_baseline_smoothing}
  \vspace{0.5cm}
\end{figure}

To achieve this, I used cross-faded concatenation of a smoothed baseline before
and after the chirp (figure \ref{fig:chirp_baseline_smoothing} B). First, I
estimated the baseline frequency EOD$f_{bl}$ by computing the median of all
values in the boundaries $b_{lower}=\min{\text{EOD}f(t)}$ and
$b_{upper}=\left|\min{\text{EOD}f(t)}\right|$. I then descended the left and
right slopes of the smoothed chirp until I either reached the baseline or the
sign of the derivative flipped, indicating that the slope goes up again. This
is not as straight forward for the end of the chirp, as a chirp often has an
undershoot. So here I descended the right slope until the sign of the
derivative flipped twice: The first time when the slope goes up again after the
undershoot and the second time when the jitter of the noisy baseline starts.
The result is a start and stop time of a chirp, which includes the undershoot
and. This can be used to extract the chirp from the centered instantaneous
frequency. To get a smoother baseline, I applied a Gaussian filter ($\sigma=60
\, \text{samples}$) to the instantaneous frequency. I then used the 120 samples
after the chirp start and the 120 samples before the chirp stop as the
cross-fade regions. In these regions, I combined the original instantaneous
frequency EOD$f(t)_{\text{orig}}$ with the smoothed instantaneous frequency
EOD$f(t)_{\text{smooth}}$ using a weighted mean with linearly increasing and
decreasing weights $w(t)$ for the cross-fade regions respectively:

\begin{equation}
  \text{EOD}f(t)_{\text{cf}} = \frac{1}{2} \left((1 - w(t)) \text{EOD}f(t)_{\text{orig}} + w(t) \text{EOD}f(t)_{\text{smooth}} \right)
\end{equation}

Finally, the baseline frequency was subtracted from EOD$f_{cf}$. The result
EOD$f_{\text{cf}}$ was now the instantaneous frequency of 9,011 chirps with a
perfectly smooth baseline that retained the natural frequency variability of a
real chirp. The remainder did not successfully pass the cross-fading and was
discarded. This happened when the descent of the slopes on either flank never
reached the baseline or a change of sign until the end of the snippet.

The resulting distribution of chirp heights and widths was severely skewed
towards small and short chirps as shown in figure
\ref{fig:chirp_overview_clustered_plot}. As I wanted the model to explicitly
learn to detect long and high chirps as well, I re-sampled the chirps to have a
uniform distribution by using a histogram equalization algorithm on both the
chirp heights and widths. This algorithm computed a 100-bin histogram for both
the distribution of chirp heights and widths. It then filled all bins by
duplicating random chirps from each bin to achieve a uniform distribution
across the bins. After histogram equalization, the dataset contained 121,380
chirps. Duplicating chirps is not ideal, as it introduces reduced variability
of large chirps. On the other hand, how a chirp appears on a spectrogram is
strongly dependent on the noise floor, how many other fish are present in the
recording and what their EOD$f$ is. So while the same chirp might appear many
times in the dataset, it will do so in many different contexts. This is much
better than having strong class imbalance. The resulting instantaneous
frequencies were then used as sources of chirps in the previously introduced
\hyperref[sec:simulatingachirpingfish]{simulation pipeline} to generate grid
recordings of multiple moving fish that produce chirps. Training data, i.e.,
spectrograms and labels, were then generated from these simulated grid
recordings as described in \hyperref[sec:spectrogramconversion]{spectrogram
conversion} and following sections. Figure \ref{fig:annotated_chirps} shows two
training samples. Each training sample is a single spectrogram image containing
many bounding boxes for chirps.

\begin{figure}
  \centering
  \includegraphics[width=\textwidth]{../figures/chirp_overview_clustered_plot.pdf}
  \caption{
    Extracted chirps for the hybrid training dataset. (A) Clustered
    instantaneous frequencies of the chirps for visualization. Chirps were
    clustered by fitting a \acrfull{GMM} by optimizing the \acrfull{BIC} on
    the first three principle components of the instantaneous frequencies. The
    upper panel shows the number of chirps per cluster. Below are the medians
    and 25\% and 75\% percentiles of the instantaneous frequencies of the
    chirps. Small and short chirps are disproportionally represented. (B) The
    distribution of chirp heights and widths as well as the troughs after the
    chirp indicated by the hue. The distribution is severely skewed towards
    small and short chirps. There is also no apparent correlation between chirp
    height and width and the trough after the chirp: Most small chirps have no
    trough, while large chirps, particularly those at around 300 \si{\hertz},
    have a trough. But chirps that are even larger or last longer have smaller
    troughs again.
  }
  \label{fig:chirp_overview_clustered_plot}
  \vspace{0.5cm}
\end{figure}

\begin{figure}
  \centering
  \includegraphics[width=\textwidth]{../figures/annotated_chirps.pdf}
  \caption{
    Two training samples from the hybrid grid datasets with real chirps. Each
    sample contains \SI{15}{\second} of the signal and captures a frequency
    range of \SI{1}{\kilo\hertz}. Bounding boxes are automatically computed by
    respecting the width and height of each underlying chirp as well as the
    \texttt{nfft} of the spectrogram to account for the artifacts introduced by
    estimating the spectrograms.
  }
  \label{fig:annotated_chirps}
  \vspace{0.5cm}
\end{figure}

\subsection{Model selection}

Deep learning and computer vision have seen a rapid development in the last
decade. This has led to a large number of different object detection models
that are tailored to different tasks and have different trade-offs in terms of
speed, accuracy and ease of use. The popular object detection models are the
R-CNN and YOLO families
\parencite{girshickRichFeatureHierarchies2014, redmonYouOnlyLook2016}. The
R-CNN family is based on the \emph{region-based convolutional neural network}
and consists of the original faster R-CNN, the R-FCN and the
Mask R-CNN. The YOLO family is based on the \emph{You Only Look Once}
model and consists of the original YOLOv1, the YOLOv2 up to
YOLOv9 and various deviations such as YOLONas and PP-YOLO,
which each have their own tweaks in model architecture and training process.

To start off, I decided to use the faster R-CNN model as it is well
documented and included in the \href{pytorch.org}{\texttt{PyTorch}} ecosystem.
The faster R-CNN model is a two-stage object detection model that first
generates a set of region proposals and then uses a region of interest pooling
layer to extract features from each region proposal. The features are then used
to predict the class and the bounding box of the object. I used the
ResNet-50 as the backbone. The features extracted by the backbone are
then used to generate the region proposals and to predict the class and the
bounding box of the object. In this case, I implemented the model training loop
using the \href{pytorch.org}{\texttt{PyTorch}} library. The model was trained
using the \texttt{Adam} optimizer with a learning rate of 0.001 and a batch
size of 8. This is the model used for the pre-annotation of the recorded
training dataset.

In addition to the faster R-CNN model, I also used the YOLOv8 and the
YOLOv9 models. These models are single-stage object detection models
that predict the class and the bounding box of the object directly from the
input image. For the YOLO models, implenting a training loop was not nessecary,
as the models were trained using the
\href{https://ultralytics.com}{\texttt{ultralytics}} library
\parencite{yolov8_ultralytics}. This includes many pre-trained models as well
as features such as a memory dependent dynamic batch sizes and a model
evaluation pipeline, which simplified the training process. All YOLOv8 as well
as the YOL0v9E models were trained and can be used as the object detection
backend in the chirp detection pipeline.

\subsection{Training \& hyperparameters}

In the final detection training dataset I combined the hybrid training dataset
with the human annotated dataset resulting in 2,159 training images. YOLOv8 and
YOLOv9 are the two main models that are used in the chirp detection pipeline,
as they are the most accurate, fast and easy to fine-tune. The models were
trained using the default hyperparameters of the \texttt{ultralytics} 8.1.23
library except for setting the image size to (800x800). Each model was trained
for 250 epochs and 5 cross-validation folds (figure
\ref{fig:yolo_crossval_training}). The batch size had to be decreased for the
larger models to fit into the memory of the GPU.

\begin{figure}
  \centering
  \includegraphics[width=\textwidth]{../figures/yolo_crossval_training.pdf}
  \caption{
    Training curves for all models that are trained for the detection step.
    Metrics were computed using 5-fold cross-validation for each model. The
    line corresponds to the median across the folds and the colored area to the
    \ordinalnum{25} and \ordinalnum{75} percentile respectively. In most case,
    the percentiles are barely visible despite training on a different subset
    of the dataset with the same random seed. Each row corresponds to a model
    and each column corresponds to a combination of loss metric and performance
    metric. The column header indicates which loss and which performance metric
    is plotted and the line colors indicate on which data (training /
    validation) the metric was evaluated. Losses decrease for all YOLOv8 models
    without too strong deviations between the training and validation loss.
    Losses and performance metrics of YOLOv9E also follow the predicted path
    but show strong fluctuations after 100 epochs are trained.
  }
  \label{fig:yolo_crossval_training}
  % \vspace{0.5cm}
\end{figure}

\subsection{Evaluation}
\label{sec:deteval}

To understand how object detection models can be evaluated, we first have to
define what a true positive and what a false positive is. In other words, how
close does a predicted bounding box have to be to the true bounding box to be
considered a true positive? And how far does a predicted bounding box have to
be from the true bounding box to be considered a false positive? To answer this
question, we can use the \acrfull{IoU} metric. The IoU between a ground truth
bounding box $y$ and a predicted bounding box $\hat{y}$ is defined as the area
of the intersection of the two bounding boxes divided by the area of the union
of the two bounding boxes:

\begin{equation}
  \text{IoU}(y, \hat{y}) = \frac{\text{area}(y \cap \hat{y})}{\text{area}(y \cup \hat{y})}
  \label{eq:iou}
\end{equation}

Traditionally, the IoU is used by thresholding it in two ways: The simple
approach is to assume that a predicted bounding box is a true positive if the
IoU is greater than a certain threshold, in most cases 0.5. A more nuanced
approach is to use a range of IoU thresholds and to compute the \acrfull{AP}
over the range of IoU thresholds, e.g., in between 0.5 and 0.95 in steps of
0.01. The average precision is defined as the area under the curve spanned by
the precision $p$ and recall $r$ as a function of the detection threshold
$\theta$, and is the most common metric to evaluate object detection models.
Precision and recall are computed from the number of true positives
($n_{\text{TP}}$), false positives ($n_{\text{FP}}$) and false negatives
($n_{\text{FN}}$):

\begin{equation}
  \begin{aligned}
    \text{p}(\theta) = \frac{n_{\text{TP}}(\theta)}{n_{\text{TP}}(\theta) + n_{\text{FP}}(\theta)} \quad \text{r}(\theta) = \frac{n_{\text{TP}}(\theta)}{n_{\text{TP}}(\theta) + n_{\text{FN}}(\theta)} 
  \end{aligned}
  \label{eq:ap}
\end{equation}

Naturally, the average precision at an IoU threshold of 0.5 (noted as
AP@.5) will result in a higher average precision than the average precision at
an IoU threshold of the previously mentioned steps (noted as
AP@[.5:.01:.95]). The latter can be formulated as the following:

\begin{equation}
  \begin{aligned}
    \text{AP@[.5:.01:.95]} = \frac{1}{N} \sum_{\theta}^{N} \text{AP}(\theta) \quad \text{where} \quad \theta \in [0.5, 0.95]
  \end{aligned}
  \label{eq:fancyap}
\end{equation}

\texttt{ultralytics} provides both metrics as the \acrfull{mAP}. This would be
the means of the average precision over all detection classes. As I am just
detecting a single class, the mAP is equivalent to the AP.

As the output of the model is a bounding box as well as the confidence for each
bounding box, I had to find the best confidence score to use as a threshold to
minimize false positives missed chirps. To do this, I used the F1 score, which
is the harmonic mean of the precision and recall.

\section{Chirp assignment}

In the assignment step, every detected bounding box needs to be attributed to a
fish in the recording. To assign each detected chirp to the identifier of the
emitter fish, I found the most informative feature to be the time dependent
power of the EOD$f$ of fish which intersect with the bounding box.

\subsection{Feature extraction}

To identify which frequency bands in chirps on spectrograms might contain the
EOD$f$ of fish, I estimated the baseline EOD$f$ of the emitter, which remains
relatively constant over short time spans. By summing the spectrogram along the
time axis to compute a power spectrum, I identified peaks representing the
EOD$f$ frequency bands of potential chirp emitters. Around each peak, a
\SI{10}{\hertz} wide band was used to calculate the average power within the
bounding box's duration for each frequency band. This feature is informative
because, during a chirp, the emitter's frequency increases, diverging from the
baseline and resulting in a power trough at the chirp's location due to the
frequency band shifting away from the baseline. However, this effect is only
partially observed due to the spectrogram's temporal resolution. While a higher
resolution would enhance this feature's utility, it would also reduce the
ability to distinguish between multiple fish due to overlapping frequency
bands. Despite this limitation, and the fact that this feature may be less
effective for detecting small chirps, it proved sufficiently reliable for
inclusion in the final pipeline. The feature extraction step is illustrated in
figure \ref{fig:assignment_feature_extraction} for a small and large chirp
respectively.

\begin{figure}
  \centering
  \includegraphics[width=\textwidth]{../figures/assignment_feature_extraction.png}
  \caption{
    Visualization of the feature processing steps starting from the predicted
    bounding boxes for a small chirp (A) and a big chirp (E) respectively. The
    power spectrum, i.e., a summation across the time axis of a bounding box,
    provides peaks at frequencies that could have been the fundamental
    frequencies of potential emitter fish (B, F). For each of these
    frequencies, the temporal evolution of the underlying power usually
    contains a trough if the current baseline frequency is that of an emitter
    fish (C, G). These power time series are then classified using a 1D
    \acrshort{CNN} (D, H). For each bounding box, the candidate frequency with
    the highest classification score is the fundamental frequency of the
    emitter fish.
  }
  \label{fig:assignment_feature_extraction}
  \vspace{0.5cm}
\end{figure}

\subsection{Training data acquisition}

The issue with this feature is that (1) deciding to which fish a chirp belongs
is not always clear on spectrograms and (2) the way this would be annotated is
not implemented in any annotation tools. So to train the assignment model, I
had to completely rely on my hybrid training dataset. To get the training data,
I computed the power of the fundamental frequency for each fish in each
bounding box during the extraction of the detection model training data. As the
chirps are simulated, the resulting dataset is truly a ground truth.

\subsection{Model selection}

To assign the chirps, I trained a simple binary classification model. For each
fish in a single bounding box, the model should return a score which indicates
the probability that a chirp belongs to the fish. Chirps are then assigned by
just accepting the fish with the highest score. As the model had to classify
simple one-dimensional arrays, I experimented with different types of
multi-layer perceptrons and 1D convolutional neural networks. The final model
is a simple CNN with two convolutional layers followed by four fully connected
layers, that are all followed by a \acrfull{ReLU} activation function except
for the output neuron, which is passed through a sigmoid. The model was build
and trained using the \href{pytorch.org}{\texttt{PyTorch}} package.

\subsection{Training \& hyperparameters}

The model was trained using the hybrid dataset. As I had much more
\enquote{non-emitters} than \enquote{emitters} in the dataset, I used a
weighted binary cross-entropy loss to balance the classes. The model was
trained with the \texttt{Adam} optimizer with a batch size of 32 and a learning
rate of 0.0001 over  11 epochs, after which the model converged (figure
\ref{fig:assignment_training} A).

\subsection{Evaluation}

To evaluate the model performance I used the \acrfull{AUC} of the \acrfull{ROC}
curve. The AUC is a metric that measures the ability of the model to
distinguish between the two classes. The AUC is calculated by plotting the
true positive rate (TPR) against the false positive rate (FPR) at various
threshold settings $\theta$ for the model confidence. The AUC is then the area
under this curve.

\begin{equation}
  \text{TPR}(\theta) = \frac{n_{\text{TP}}(\theta)}{n_{\text{TP}}(\theta) + n_{\text{FN}}(\theta)} \quad \text{FPR}(\theta) = \frac{n_{\text{FP}}(\theta)}{n_{\text{FP}}(\theta) + n_{\text{TN}}(\theta)} \label{eq:roc}
\label{eq:aucroc}
\end{equation}

The AUC ranges from 0 to 1, where 0.5 corresponds to a model that randomly
assigns the classes and 1 corresponds to a model that always assigns the
correct class. The AUC is a good metric to evaluate the performance of the
model as it is independent of the threshold setting and thus can also be used
to choose the best threshold setting for the model in production. As the AUC
can provide overly optimistic performance estimates for strongly imbalanced
classes, I also used the average precision, i.e., the area under the
precision-recall curve as well as the F1 scores (as explained in section
\ref{sec:deteval}) to evaluate the models performance and find the best
threshold.

\section{Behavioral analyses}

All behavioral analyses on the Raab dataset were conducted as described in
\textcite{raabElectrocommunicationSignalsIndicate2021}. The agonistic
interactions were annotated using
\href{https://github.com/olivierfriard/BORIS}{\texttt{BORIS}}, a free and
open-source software designed to facilitate the coding of behaviors in animal
behavior studies. The EOD$f$s of the fish were tracked using the
\href{https://github.com/tillraab/wavetracker}{\texttt{wavetracker}} package
for \texttt{Python}. Annotated behaviors and tracked EOD$f$s were already
available. I detected and assigned chirps to individual fish using the here
introduced
\href{https://github.com/weygoldt/chirpdetector/}{\texttt{chirpdetector}}
package. The chirp rates during the time of agonistic events were computed by
convolving the chirps of each fish with a Gaussian kernel with a standard
deviation of one second and a height of one. The \ordinalnum{1} and
\ordinalnum{99} confidence intervals were estimated by jackknife resampling
corresponding of an $\alpha$-level of 1\%. The null hypothesis, which is that
the chirp rates are independent of the on- and offset of agonistic events was
tested using bootstrapping. These methods are introduced in detail in
\textcite{raabElectrocommunicationSignalsIndicate2021}. Statistical analyses of
chirp counts and chasing durations were performed using \acrfullpl{GLMM}
implemented in the
\href{https://cran.r-project.org/web/packages/lme4/index.html}{\texttt{lme4}}
package in \texttt{R v4.3.3} with the pairing ID as a random factor
\parencite{batesFittingLinearMixedEffects2015}. To test whether the competition
status of two fish had an effect on the number of chirps each fish produced, I
used a GLMM of the Poisson family with a logarithmic link function. To model
the effect the duration of chasing had on whether loser produced chirps or not,
I weighted each data point (duration of a single chasing and the presence of
chirps) by (1) how many chasing events were produced by each pairing and (2)
the fraction of chasings in which chirps were produced, to not bias the model
to a specific trial. I then used a Binomial GLMM with a Logit link function.
The effect size is reported as Cohen's $d$. 

